\section{Introduction}

The random switches in a shuffle do more than move individual cards. They define a bijection between all starting and ending positions. If the paths of some cards are revealed, the other cards still undergo a random bijection, but its law depends on the revealed paths. We study when that entire conditional bijection becomes nearly uniform. This asks for information that is absent from a one-card mixing estimate: all remaining labels must be randomized together, after the surrounding motion has been observed.

The Thorp shuffle makes this distinction concrete. Split a deck into two equal halves, pair the cards in corresponding positions, and independently choose the order of each pair before interleaving them. For $n=2^d$ positions, undoing a deterministic rotation of the binary coordinates identifies successive shuffles with fair switches in the cyclic directions of the binary cube. A \emph{sweep} visits all $d$ directions once and costs $d$ physical shuffles. Each individual card is uniform after one sweep. The dependence among the routes of different cards remains substantial.

Write $U_m$ for uniform measure on the symmetric group $S_m$, and use total variation with its factor $1/2$. The mixing time $t_{\mathrm{mix}}(d)$ is the least number of physical shuffles at which the full-deck law is within $1/4$ of $U_n$ from every deterministic start. Fix $m=n/2$ outside labels. Given their full trajectories $\mathcal E$ through a time interval, list the unoccupied slots at each endpoint in increasing binary order. The switches map the initial free slots bijectively onto the final free slots; the resulting permutation of their indices has a conditional law $\nu_{\mathcal E}$ on $S_m$. Its definition uses the slot indices, independently of the order of the labels entering the interval. Section~\ref{sec:conditional} proves both this construction and the composition rule on consecutive intervals.

The main contribution is the information retained after the outside paths are revealed. An independent companion~\cite{optimal-order} gives full-deck mixing after $1600d$ physical shuffles. That conclusion concerns the unconditioned deck law; the theorem here controls the entire remaining assignment conditional on a path environment, with its own explicit time bound.

\begin{theorem}[Conditional permutation and full-deck mixing]\label{thm:methods}
Let $n=2^d$, and start the shuffle from any deterministic deck.
After $J(s_0+1)$ sweeps, where $s_0=400$ and $J=40000$, the conditional free-slot law defined above satisfies
\[
 \sup_{\text{initial layouts}}\E_{\mathcal E}
       \|\nu_{\mathcal E}-U_{n/2}\|_{\TV}\longrightarrow0
       \qquad(d\longrightarrow\infty).
\]
Consequently the full-deck law approaches uniform, uniformly in the starting deck, after $16\,040\,400d$ physical shuffles. For every integer $t\ge0$, after $t$ physical shuffles the full-deck distance from uniform is at least
$1-2^{tn/2}/n!$. Thus the full-deck mixing time is of order $d$.
\end{theorem}

The conditional conclusion says that the free-slot permutation is nearly independent of the revealed environment. Indeed, if $g$ is that permutation, then
\[
 \left\|\mathcal L(\mathcal E,g)
       -\mathcal L(\mathcal E)\otimes U_{n/2}\right\|_{\TV}
 =\E_{\mathcal E}\|\nu_{\mathcal E}-U_{n/2}\|_{\TV}.
\]
The equality follows by summing first over the environment and then over the permutation. For each fixed $\varepsilon>0$, Markov's inequality also shows that the probability of a conditional distance exceeding $\varepsilon$ tends to zero, uniformly in the initial layout. Rare path environments can nevertheless leave too few useful switches. The proof handles them by deleting a small amount of probability, with the loss measured under the actual outside-path law; it does not assume that a typical free set stays typical after further conditioning.

The first proof has three steps. We first mix any ordered half-deck in $400$ sweeps from every start. A two-copy calculation for a single sweep gives the exact collision kernel of a tagged free card. Splitting it into averages within large subcubes and fluctuations inside those subcubes produces a two-sweep contraction. Next, we alternate a preparation interval of $400$ sweeps with one active sweep. Truncating the preparation law gives domination by a uniform free set and an independent uniform slot permutation. The active sweep forces many coin agreements whenever two copies have many common endpoints. A symmetric-group character bound converts this into a normalized trace contraction. The preparation domination then turns that trace estimate into a matrix contraction which can be iterated while the outside paths continue to be revealed. Finally, we couple the outside half-decks and use the conditional result to couple the remaining labels.

The three subsequent constructions change a specific part of this mechanism. Their separate quantitative statements and time bounds are given with their proofs. Reset minorization uses an exact uniform component in a large marginal to create central relative-permutation increments. Its conditional collision identity gives uniformity even after all reset coins are revealed. Trajectory overlays preserve the true shuffle law by adding predictable independent switches to a base trajectory system. They first randomize one group of indices and then use three rounds on transverse partitions; a local limit estimate for contingency tables completes that route. The last construction conditions on the unlabelled occupancy history of two colors. Given that history, the two internal half-permutations are independent. Conditional list estimates and a signed-tabloid transfer then control a two-block Fourier matrix. These are different conditional objects, and each construction includes its own mass and independence argument.

The shuffle originates in Thorp's study of shuffling and Faro~\cite{thorp}. Morris proved an $O(d^{44})$ bound for power-of-two decks~\cite{morris44}; Montenegro and Tetali gave the intermediate $O(d^{29})$ bound~\cite[Theorems~6.14--6.15]{montenegro-tetali2006}. Morris subsequently obtained an $O((\log n)^4)$ bound for every even deck size~\cite{morris4} and an $O(d^3)$ bound for $n=2^d$~\cite{morris3}.

For partial permutations, Czumaj and V\"ocking proved that the endpoint positions of any prescribed $\lfloor cn\rfloor$ cards mix in $O(d^2)$ physical Thorp shuffles, for each fixed $0<c<1$ and $n=2^d$~\cite{czumaj-vocking2014}. Morris's chameleon analysis already studies a tagged card conditional on an occupancy history in an auxiliary zigzag shuffle~\cite[preprint, Section~4, Lemma~3]{morris44}. Partial permutation laws also appear in the enciphering work of Morris, Rogaway and Stegers~\cite[Theorem~1 and Lemma~2, proved in Appendix~A]{mrs}. Hoang, Morris and Rogaway use conditional squared energies and sequential total-variation comparison in their analysis of swap-or-not~\cite[Section~3]{hmr2014}. Their random-key chain differs from the deterministic coordinate schedule here. Our conditional permutation estimates require the additional bijection, truncation and character arguments developed below.

The character input is the uniform Larsen--Shalev estimate~\cite[Theorem~1.1]{larsen-shalev2008}, in the precise fixed-slack formulation recorded by Keller, Lifshitz and Sheinfeld~\cite[Definition~2.1 and Theorem~2.2]{keller-lifshitz-sheinfeld2024}. Standard branching, Schur orthogonality and hook-length formulas are used in the conventions of~\cite{sagan,etingof}. Lemmas~\ref{lem:character-fixed-points} and~\ref{lem:degree-sum} collect the character consequences and an elementary reciprocal-square degree estimate used across the constructions. The latter is a special case of the general summability theorem of Liebeck and Shalev~\cite[Theorem~2.6]{liebeck-shalev2004}.

Sections~\ref{sec:conditional}--\ref{sec:burn-completion} give the entire first route. Sections~\ref{special:reset} and~\ref{special:overlay} give the reset and overlay constructions. Section~\ref{sec:two-halves} identifies precisely the conditional path theorem and the abstract signed-tabloid transfer used in the two-half construction. The companion on conditional coordinate sweeps~\cite{conditional-sweeps} studies prescribed trajectories in a different, near-uniform line model. It retains a falling-factorial potential through the conditional moment induction; with no prescribed paths, the resulting dimension estimate is then transferred analytically to binary sweeps. Its theorem is not an input to the present kernels.
